% =====================================================================
\chapter{Phương pháp đo và tiêu chí nghiệm thu}
\label{ch:eval}
% =====================================================================

Chương này cố định cách đo hai kết quả chính. Mọi tham số đánh dấu \TBD{} phải
được chốt và khoá \emph{trước} khi chạy trên tập held-out (mục~\ref{sec:prereg}).

% ---------------------------------------------------------------------
\section{KPI-1 · Recall lỗi trong top 20\% frame}
\label{sec:kpi1}

\subsection{Định nghĩa}
Với snapshot held-out có \(N\) frame, tập lỗi reference \(\mathcal{E}\) (chỉ E1–E3,
đã loại ignore theo BR-07), ranking \(\pi\) do score version đã khoá sinh ra:
\[
  n_{20} = \lceil 0{,}2\,N \rceil, \qquad
  \mathrm{Recall@20\%} = \frac{\bigl|\{e \in \mathcal{E} : \mathrm{rank}_\pi(\phi(e)) \le n_{20}\}\bigr|}{|\mathcal{E}|}.
\]

\begin{itemize}
  \item \textbf{Đơn vị đếm:} lỗi (BR-01), không phải frame.
  \item \textbf{Làm tròn:} lấy trần \(\lceil 0{,}2N \rceil\).
  \item \textbf{Hoà điểm:} phá hoà tất định (FR-RNK-04); không chọn thứ tự có lợi.
  \item \textbf{Điều kiện tính:} \(|\mathcal{E}| \ge E_{min}\) (\TBD[-K4]); dưới
        ngưỡng ghi ``không đủ mẫu''.
  \item \textbf{Ý nghĩa:} KPI-1 đo \emph{khả năng xếp hạng} của công cụ, độc lập
        với việc reviewer có thực sự tìm thấy lỗi. Hiệu quả phát hiện thực tế đo
        bằng KPI-2b trong thí nghiệm effort.
\end{itemize}

\subsection{Đối chứng}
\begin{xltabular}{\linewidth}{@{}L{3.4cm}Y@{}}
\caption{Ranking đối chứng cho KPI-1}\label{tab:baselines1}\\
\toprule
\textbf{Đối chứng} & \textbf{Cách tạo và ý nghĩa}\\
\midrule
\endfirsthead
\toprule \textbf{Đối chứng} & \textbf{Cách tạo}\\ \midrule
\endhead
\bottomrule
\endfoot
Ngẫu nhiên & Hoán vị ngẫu nhiên đều, \(S \ge 1000\) seed. Kỳ vọng xấp xỉ \(n_{20}/N \approx 20\%\); báo cáo trung bình và phân vị 2{,}5\%–97{,}5\%, không dùng con số 20\% cố định.\\
Heuristic mật độ & Sắp theo số annotation trong frame giảm dần. Kiểm tra công cụ có tốt hơn quy tắc ``frame đông thì nhiều lỗi''.\\
Heuristic Detector & Sắp theo confidence lớn nhất của candidate E1/E2 trong frame (không cộng điểm nền, khác \code{score\_v0}). Kiểm tra giá trị của hiệu chỉnh và gộp nhiều nguồn.\\
Oracle (tham khảo) & Sắp theo số lỗi thật trong frame; cho biết cận trên Recall@20\% có thể đạt.\\
\end{xltabular}

\subsection{Khoảng tin cậy và tiêu chí đạt}
\begin{itemize}
  \item Khoảng tin cậy 95\% bằng \emph{cluster bootstrap} theo video nguồn
        (\(B \ge 1000\), seed khoá trước), vì các frame cùng video tương quan.
        Mỗi lần lặp: lấy lại có hoàn lại toàn bộ video; frame của video lặp lại
        được nhân bản với khoá phá hoà riêng; tính lại \(N\), \(n_{20}\), sắp lại
        theo \(s(f)\) đã có (không học lại điểm) và tính lại các ranking đối chứng
        trên mẫu đó. Recall của ranking ngẫu nhiên lấy theo kỳ vọng \(n_{20}/N\)
        của từng mẫu.
  \item Tính hiệu số \(\Delta = \mathrm{Recall@20\%}_{\pi} - \mathrm{Recall@20\%}_{\text{đối chứng}}\)
        trên cùng mẫu bootstrap.
  \item \textbf{Đạt KPI-1} khi đồng thời: (i) cận dưới khoảng tin cậy của
        Recall@20\% $\ge$ ngưỡng \TBD[-K1]; (ii) cận dưới khoảng tin cậy của
        \(\Delta\) so với ngẫu nhiên và với heuristic mật độ đều \(> 0\).
  \item Báo cáo thêm theo nhóm lỗi E1/E2/E3 và theo slice; nhóm có ít hơn
        \(E_{min,g}\) lỗi chỉ báo cáo, không kết luận.
\end{itemize}

\begin{figure}[H]
\centering
\begin{tikzpicture}
\begin{axis}[
  width=0.82\linewidth, height=6.2cm,
  xlabel={\sffamily\small Tỉ lệ frame đã review theo ranking (\%)},
  ylabel={\sffamily\small Tỉ lệ lỗi reference nằm trong phần đã review (\%)},
  xmin=0, xmax=100, ymin=0, ymax=100,
  xtick={0,20,40,60,80,100}, ytick={0,20,40,60,80,100},
  grid=major, grid style={lxline},
  legend style={font=\sffamily\scriptsize, at={(0.98,0.03)}, anchor=south east},
  tick label style={font=\sffamily\scriptsize},
]
\addplot[lxgray, thick, dashed] coordinates {(0,0) (100,100)};
\addlegendentry{Ngẫu nhiên (kỳ vọng)}
\addplot[lxamber, thick] coordinates {(0,0) (10,17) (20,31) (40,52) (60,70) (80,86) (100,100)};
\addlegendentry{Heuristic mật độ (minh hoạ)}
\addplot[lxblue, very thick] coordinates {(0,0) (5,24) (10,40) (20,62) (40,82) (60,92) (80,97) (100,100)};
\addlegendentry{Ranking M13 (minh hoạ)}
\addplot[lxgreen, thin] coordinates {(0,0) (5,48) (10,78) (15,95) (20,100) (100,100)};
\addlegendentry{Oracle (minh hoạ)}
\draw[lxred, dashed] (axis cs:20,0) -- (axis cs:20,100);
\node[font=\sffamily\scriptsize, text=lxred, anchor=west] at (axis cs:21,8) {$k = 20\%$};
\end{axis}
\end{tikzpicture}
\caption{Đường Recall@k và cách đọc KPI-1 (số liệu \emph{minh hoạ}, không phải kết quả đo)}
\label{fig:recallcurve}
\end{figure}

% ---------------------------------------------------------------------
\section{KPI-2 · Giảm công sức review với chất lượng tương đương}
\label{sec:kpi2}

\subsection{Định nghĩa effort}
Effort là \emph{thời gian thao tác hoạt động} ghi tự động (FR-EVL-06), loại các
khoảng không thao tác dài hơn \(t_{idle}\). Thời gian được gán vào đúng một hoạt
động để không cộng hai lần.

\begin{xltabular}{\linewidth}{@{}L{1.4cm}L{4.2cm}Y C{1.6cm}@{}}
\caption{Thành phần effort}\label{tab:effort}\\
\toprule
\textbf{Ký hiệu} & \textbf{Hoạt động} & \textbf{Ghi chú} & \textbf{Tính vào \(T\)}\\
\midrule
\endfirsthead
\toprule \textbf{Ký hiệu} & \textbf{Hoạt động} & \textbf{Ghi chú} & \textbf{Tính vào \(T\)}\\ \midrule
\endhead
\bottomrule
\endfoot
\(t_{view}\) & Xem frame, không gắn issue cụ thể & Gồm cả frame sạch. & Có\\
\(t_{tp}\) & Xử lý issue cuối cùng là lỗi thật & Từ khi mở issue đến khi lưu quyết định. & Có\\
\(t_{fa}\) & Xử lý cảnh báo sai & Issue bị bác bỏ hoặc phân xử là không lỗi. Là một phần của thời gian review, tách ra để báo cáo. & Có\\
\(t_{adj}\) & Phân xử của QA Lead & Cho các issue chuyển cấp. & Có\\
\(t_{rc}\) & Kiểm lại sau sửa & Reviewer xem kết quả re-check và verify. & Có\\
\(t_{fix}\) & Annotator sửa trên CVAT & Ghi bằng thời điểm mở deep link đến khi bấm \emph{Đã sửa}; báo cáo riêng và trong phân tích độ nhạy. & Không\\
\end{xltabular}

\[
  T = \sum \left(t_{view} + t_{tp} + t_{fa} + t_{adj} + t_{rc}\right), \qquad
  \text{KPI-2} = 1 - \frac{T_A}{T_B}.
\]

Chỉ số phụ: lỗi xác nhận trên giờ review, \(t_{fa}/T\) (tỉ trọng thời gian cho cảnh
báo sai), KPI-2b (recall phát hiện thực tế).

\subsection{Thiết kế thí nghiệm}
Thí nghiệm so sánh \emph{hai quy trình} (baseline và assisted) trên cùng dữ liệu
và đo chất lượng cuối trên toàn tập. Lợi ích đo được là của cả quy trình
(ranking, evidence, guideline, điểm dừng), không quy riêng cho ranking.

\begin{figure}[H]
\centering
\begin{tikzpicture}[node distance=4mm and 8mm]
  \node[boxg, text width=26mm] (d) {Dữ liệu thí nghiệm\\\scriptsize chia ngẫu nhiên D1, D2\\phân tầng theo video, slice};
  \node[boxb, text width=30mm, right=12mm of d, yshift=10mm] (g1p1) {\textbf{Chuỗi AB · đợt 1}\\Baseline trên D1};
  \node[boxgr, text width=30mm, right=of g1p1] (g1p2) {\textbf{Chuỗi AB · đợt 2}\\Assisted trên D2};
  \node[boxgr, text width=30mm, right=12mm of d, yshift=-10mm] (g2p1) {\textbf{Chuỗi BA · đợt 1}\\Assisted trên D1};
  \node[boxb, text width=30mm, right=of g2p1] (g2p2) {\textbf{Chuỗi BA · đợt 2}\\Baseline trên D2};
  \draw[arr] (d.east) -- (g1p1.west);
  \draw[arr] (d.east) -- (g2p1.west);
  \draw[arr] (g1p1) -- (g1p2);
  \draw[arr] (g2p1) -- (g2p2);
  \node[font=\sffamily\scriptsize, align=center, below=3mm of g2p2, xshift=-20mm, text=lxgray]
    {Reviewer được gán ngẫu nhiên vào chuỗi AB hoặc BA (phân tầng theo kinh nghiệm).\\Mỗi reviewer gặp mỗi tập dữ liệu và mỗi điều kiện đúng một lần.};
\end{tikzpicture}
\caption{Thiết kế chéo (crossover) cho thí nghiệm effort}
\label{fig:crossover}
\end{figure}

\begin{xltabular}{\linewidth}{@{}L{1.3cm}Y@{}}
\caption{Quy tắc thí nghiệm effort}\label{tab:exprules}\\
\toprule
\textbf{Mã} & \textbf{Quy tắc}\\
\midrule
\endfirsthead
\toprule \textbf{Mã} & \textbf{Quy tắc}\\ \midrule
\endhead
\bottomrule
\endfoot
EX-01 & Reviewer được gán \emph{ngẫu nhiên} (seed khoá trước) vào chuỗi AB hoặc BA, phân tầng theo kinh nghiệm và điểm calibration. AS-05 (4 reviewer) là mức sàn; số reviewer thực tế lấy từ power analysis (TBD-12).\\
EX-02 & D1, D2 chia ngẫu nhiên theo video nguồn, phân tầng theo slice và số lỗi reference; không giao với tập hiệu chỉnh. Trong mỗi tập, frame được chia cho reviewer cùng nhánh theo khối video ngẫu nhiên, mỗi frame đúng một reviewer.\\
EX-03 & Mỗi nhánh làm trên \emph{bản sao annotation riêng} (CVAT task riêng) cùng xuất phát từ snapshot đầu vào.\\
EX-04 & \textbf{Baseline:} review toàn bộ frame theo thứ tự gốc, Workspace ẩn ranking, risk score và evidence engine; quy trình sửa và verify như assisted.\\
EX-05 & \textbf{Assisted:} review tập \(R = T_{k^\ast}(\pi) \cup U\), với \(U\) là \(\lceil rN \rceil\) frame chọn ngẫu nhiên đều (seed khoá trước) từ các frame \emph{ngoài} \(T_{k^\ast}(\pi)\). Thứ tự review: hết \(T_{k^\ast}\) theo rank, rồi tới \(U\).\\
EX-06 & \textbf{Chọn \(k^\ast\):} trên tập hiệu chỉnh, với lưới \(k \in \{10\%, 15\%, \dots, 100\%\}\), \(k^\ast\) là giá trị nhỏ nhất mà residual dự kiến của assisted (ngoài fold) không vượt residual baseline quá \(\delta/2\). \(k^\ast\), \(r\), seed khoá trước thí nghiệm.\\
EX-07 & \textbf{Điểm kết thúc đo:} một nhánh kết thúc khi mọi frame trong phạm vi đã review xong, mọi issue (kể cả issue mới từ re-check và từ QC run cuối) đã có kết luận, và mọi rework đã verify. Effort tính đến thời điểm đó.\\
EX-08 & \textbf{Chống carryover:} khoá guideline và kho case ở một version trong suốt thí nghiệm; case phân xử phát sinh trong thí nghiệm không được hiển thị cho nhánh khác đến khi kết thúc; reviewer không trao đổi về dữ liệu thí nghiệm; có buổi làm quen như nhau trước đợt 1.\\
EX-09 & \textbf{Nguồn lực như nhau:} cùng QA Lead phân xử, cùng quy tắc verify, cùng quyền tra guideline ở hai nhánh.\\
EX-10 & Reviewer không biết reference; người lập reference không tham gia thí nghiệm trên cùng dữ liệu.\\
\end{xltabular}

\subsection{Phân tích effort}
\begin{itemize}
  \item \textbf{Estimand chính:} tỉ số \(T_A/T_B\) của tổng effort trên cùng khối
        lượng dữ liệu; \(\text{KPI-2} = 1 - T_A/T_B\).
  \item \textbf{Mô hình:} hồi quy hỗn hợp trên \(\log T_{r,p}\) (effort của
        reviewer \(r\) ở đợt \(p\), chuẩn hoá theo số frame của tập):
        \(\log T_{r,p} = \mu + \tau\cdot\mathrm{arm} + \gamma\cdot\mathrm{period}
        + \eta\cdot\mathrm{dataset} + u_r + \varepsilon\), với \(u_r\) là hiệu ứng
        ngẫu nhiên của reviewer. \(\exp(\tau)\) ước lượng \(T_A/T_B\).
  \item \textbf{Khoảng tin cậy:} từ mô hình trên, kiểm tra lại bằng bootstrap hai
        tầng (lấy lại reviewer trong từng chuỗi, rồi lấy lại khối video trong
        từng tập), giữ nguyên cặp hai điều kiện của mỗi reviewer.
  \item \textbf{Phân tích độ nhạy định trước:} (i) chỉ dùng đợt 1 (so sánh song
        song, không có carryover); (ii) hai giá trị \(t_{idle}\); (iii) cộng
        \(t_{fix}\) vào \(T\).
\end{itemize}

\subsection{Chất lượng tương đương}
Trong thí nghiệm, reference phủ toàn bộ dữ liệu nên residual được đo \emph{trực
tiếp trên toàn tập}, không ước lượng từ mẫu. Sau khi mỗi nhánh kết thúc, so
annotation cuối của nhánh với GT theo cùng quy tắc ở mục~\ref{sec:counting}:
\[
  \rho = \frac{\text{số lỗi E1–E3 còn lại sau quy trình}}{\text{số đối tượng GT không ignore}} .
\]
Lỗi còn lại gồm lỗi ban đầu chưa sửa và lỗi mới phát sinh do sửa. \(\rho\) có thể
vượt 100\% trong trường hợp cực đoan (nhiều box trùng); mẫu số bằng 0 thì không
tính. \(\rho\) khác G-4: G-4 là guardrail vận hành ước lượng từ lát kiểm tra
ngẫu nhiên khi không có reference toàn tập.

\begin{itemize}
  \item \textbf{Kiểm định non-inferiority một phía:}
        \(H_0: \rho_A - \rho_B \ge \delta\) với \(H_1: \rho_A - \rho_B < \delta\),
        mức \(\alpha = 0{,}025\). Tương đương: bác \(H_0\) khi cận trên khoảng
        tin cậy 95\% hai phía của \(\rho_A - \rho_B\) nhỏ hơn \(\delta\).
  \item \textbf{Khoảng tin cậy:} bootstrap theo khối video trong từng tập, giữ
        cặp hai nhánh trên cùng dữ liệu.
  \item \textbf{Chọn \(\delta\):} Product Owner và QA Lead chốt dựa trên mức
        residual chấp nhận được về nghiệp vụ (ví dụ một tỉ lệ của residual
        baseline đo ở pilot nhỏ), trước thí nghiệm (TBD-K3).
  \item \textbf{G-2b — lỗi nghiêm trọng:} số lỗi nghiêm trọng của reference còn
        lại trong annotation cuối (kể cả lỗi chưa từng có issue) ở nhánh assisted
        không lớn hơn nhánh baseline. G-2b khác G-2 (gate: issue nghiêm trọng đã
        xác nhận còn mở = 0).
  \item \textbf{G-5 — box thừa:} số box thừa (BR-06) trong annotation cuối, kể cả
        box phát sinh do sửa, báo cáo theo nhánh; nhánh assisted không được nhiều
        hơn baseline quá \(\delta_{fp}\) (TBD-K3).
  \item Ngưỡng residual \(\le 5\%\) và đồng thuận \(\ge 90\%\) là guardrail
        khởi điểm; chúng không thay thế kiểm định trên.
\end{itemize}

\subsection{Tiêu chí đạt KPI-2}
\textbf{Đạt KPI-2} khi đồng thời: (i) bác được \(H_0\) non-inferiority về
\(\rho\); (ii) cận dưới khoảng tin cậy 95\% của \(1 - T_A/T_B\) $\ge$ ngưỡng
\TBD[-K2]; (iii) đạt G-2b và G-5; (iv) kết luận không đổi chiều ở phân tích chỉ
dùng đợt 1.

\begin{notebox}[Cỡ mẫu]
Số frame mỗi tập và số reviewer được tính từ phương sai effort và residual đo ở
pilot nhỏ trên tập hiệu chỉnh (power analysis 80\%, có tính hệ số thiết kế do
phân cụm theo video và reviewer). Cỡ mẫu chốt trước thí nghiệm (\TBD[-12]).
\end{notebox}

\section{Đăng ký trước (pre-registration)}
\label{sec:prereg}
Các mục sau phải được khoá version trong LabelX \emph{trước} khi chạy trên tập
held-out hoặc bắt đầu thí nghiệm. Thay đổi sau khoá phải ghi lý do và kết quả
dùng phiên bản cũ vẫn được báo cáo.

\begin{xltabular}{\linewidth}{@{}L{5.2cm}Y L{2.4cm}@{}}
\caption{Danh sách khoá trước khi đo}\label{tab:prereg}\\
\toprule
\textbf{Mục} & \textbf{Nội dung} & \textbf{Người khoá}\\
\midrule
\endfirsthead
\toprule \textbf{Mục} & \textbf{Nội dung} & \textbf{Người khoá}\\ \midrule
\endhead
\bottomrule
\endfoot
Detector & Artifact, checksum, mapping lớp, \(\tau_{conf}\), xác nhận không huấn luyện trên held-out. & Data/Model Owner\\
Score version & Đặc trưng, tham số \(\mathbf{w}_g, b_g\), mô hình điểm nền \(h(f)\), \(\tau_{E1}, \tau_{E2}\), seed phá hoà. & QC Admin\\
Reference & GT version, mapping, \(\tau_m\), \(a_{min}\), \(E_{min}\). & QA Lead\\
Ngưỡng KPI & K1, K2, \(\delta\). & Product Owner\\
Thí nghiệm & Seed chia D1/D2 và gán chuỗi AB/BA, \(k^\ast\), \(r\), \(t_{idle}\), cỡ mẫu, mô hình phân tích, phân tích độ nhạy, \(\delta\), \(\delta_{fp}\). & Product Owner\\
\end{xltabular}

% ---------------------------------------------------------------------
\section{Ma trận nghiệm thu}
\label{sec:acceptance}
\begin{xltabular}{\linewidth}{@{}L{1.2cm}Y L{1.6cm}L{3cm}@{}}
\caption{Tiêu chí nghiệm thu}\label{tab:acceptance}\\
\toprule
\textbf{Mã} & \textbf{Tiêu chí kiểm chứng được} & \textbf{Mục tiêu} & \textbf{Yêu cầu}\\
\midrule
\endfirsthead
\toprule \textbf{Mã} & \textbf{Tiêu chí} & \textbf{Mục tiêu} & \textbf{Yêu cầu}\\ \midrule
\endhead
\bottomrule
\endfoot
AC-01 & Tạo hai lần QC Run cùng snapshot, config, score version, seed cho ra candidate, issue, ranking giống hệt: so hash của nội dung đã chuẩn hoá (bỏ ID sinh tự động, thời điểm, run\_id; sắp theo khoá tự nhiên). & R-01 & FR-AGG-02, FR-RNK-03\\
AC-02 & Ép lỗi một shard rồi retry: số issue không đổi, không có bản trùng. & R-01 & FR-AGG-02\\
AC-03 & Sửa annotation trên CVAT trong lúc tạo snapshot: snapshot không được khoá, báo job drift. & R-01 & FR-SNP-04\\
AC-04 & Mọi candidate hiển thị trong Workspace có evidence; không issue nào được đóng hoặc xác nhận mà không có quyết định của người. & R-02 & FR-REV-07, FR-REV-08\\
AC-05 & Reviewer là assignee của job tại snapshot gửi quyết định qua API: bị từ chối, có audit. & R-05 & FR-SEC-03\\
AC-06 & Issue chỉ Đã đóng sau verify trên revision mới; báo cáo trỏ QC run cuối. & R-03 & FR-RWK-05, FR-RWK-06\\
AC-07 & Engine Not checked/Failed hiển thị đúng trạng thái; gate không đạt khi thiếu dữ liệu. & R-06 & FR-AGG-05, FR-GTE-02\\
AC-08 & Recall@20\% trên held-out đạt tiêu chí mục~\ref{sec:kpi1}, kèm đối chứng và khoảng tin cậy. Chỉ đánh giá được sau khi TBD-K1, TBD-K4 đã chốt; trước đó trạng thái là ``chưa xác định''. & KPI-1, R-04 & FR-EVL-07…10\\
AC-09 & Thí nghiệm effort đạt tiêu chí mục~\ref{sec:kpi2}. Chỉ đánh giá được sau khi TBD-K2, TBD-K3, TBD-12 đã chốt; trước đó trạng thái là ``chưa xác định''. & KPI-2, R-07 & FR-EVL-11…13\\
AC-10 & Báo cáo hiệu quả có đủ cỡ mẫu, mẫu số, phương pháp, provenance; random và risk báo cáo riêng. & R-04 & FR-RPT-01…04\\
AC-11 & Audit log có đủ actor, đối tượng, trước/sau, lý do cho mọi quyết định, phân xử, waiver, khoá reference. & R-05 & FR-SEC-05\\
\end{xltabular}
